\chapter{Introdução} % ELEMENTO OBRIGATÓRIO

A análise histopatológica é o padrão-ouro para o diagnóstico de diversas doenças, em especial do câncer. Entretanto, por ser realizada de forma predominantemente manual, é um processo demorado, dependente da experiência do especialista e sujeito a divergências de interpretação. Diante desse problema, sistemas de diagnóstico auxiliado por computador (CAD, do inglês \textit{Computer-Aided Diagnosis}) vêm sendo desenvolvidos para oferecer uma segunda opinião ao patologista, o que confere relevância a esta área de pesquisa ao tornar o diagnóstico mais ágil e menos subjetivo.

Nesse contexto, os modelos de aprendizado profundo se destacam, mas exigem grandes volumes de dados, raramente disponíveis na área médica. Por outro lado, os atributos fractais descrevem de forma compacta a complexidade e a organização dos tecidos, porém resultam em vetores numéricos que não aproveitam a capacidade desses modelos de reconhecer padrões em imagens. Justifica-se, portanto, a conversão desses vetores em representações visuais como forma de unir as duas abordagens. Para essa tarefa, o Campo de Transição de Markov (MTF, do inglês \textit{Markov Transition Field}) mostra-se promissor e, até o momento, ainda não foi explorado em associação a atributos fractais.

Assim, o objetivo deste trabalho é investigar a classificação de imagens histológicas por meio da associação entre Geometria Fractal e o MTF, comparando-a a outras formas de representação e a diferentes arquiteturas de aprendizado profundo. Desse modo, busca-se contribuir para sistemas de apoio ao diagnóstico mais precisos e menos dependentes de grandes bases de dados. Por fim, o estudo limita-se a bases públicas de tamanho reduzido e não contempla a validação clínica dos resultados.
